% Source proof digest supplied by the user: 
% Historical journal record: 
\documentclass{article}

\usepackage{amsmath,amsthm,amssymb}

\newtheorem{theorem}{Theorem}[section]
\newtheorem{proposition}[theorem]{Proposition}
\newtheorem{lemma}[theorem]{Lemma}
\newtheorem{corollary}[theorem]{Corollary}
\theoremstyle{remark}
\newtheorem{remark}[theorem]{Remark}

\newcommand{\ind}{\mathbf{1}}
\newcommand{\Ccal}{\mathcal{C}}
\newcommand{\Ocal}{\mathcal{O}}

\title{Odd Cycles in the Coprime Graph at the Sharp Threshold}
\author{}
\date{}

\begin{document}

\maketitle

\begin{abstract}
For \(A\subseteq\{1,\ldots,n\}\), let \(G(A)\) be the graph in which two
integers are adjacent when they are coprime. We prove that if
\[
 |A|>
 \left\lfloor\frac n2\right\rfloor+
 \left\lfloor\frac n3\right\rfloor-
 \left\lfloor\frac n6\right\rfloor,
\]
then \(G(A)\) contains every odd cycle of length at most
\[
2\left(
 \left\lfloor\frac n3\right\rfloor-
 \left\lfloor\frac n6\right\rfloor
\right)+1.
\]
Both the cardinality threshold and this exact ceiling are best possible.
When \(6\mid n\), the largest forced length is \(n/3+1\), resolving the
Erdős--Sárközy conjecture with the sharp constant. The proof constructs
one covering path whose prefixes give all required cycles. Its main
ingredients are a Hall-type connector lemma, an upper-tail mean estimate
for a multiplicative function, a selection lemma based on majorization,
and a two-reservoir transportation argument. Two numerical inequalities
are established by finite exact-rational inclusion--exclusion
computations.
\end{abstract}

\section{Introduction}

For a positive integer \(n\), write \([n]=\{1,\ldots,n\}\). Given
\(A\subseteq[n]\), its \emph{coprime graph} \(G(A)\) has vertex set \(A\),
with two distinct vertices adjacent if and only if they are coprime. Put
\[
 T(n)=
 \left\lfloor\frac n2\right\rfloor+
 \left\lfloor\frac n3\right\rfloor-
 \left\lfloor\frac n6\right\rfloor
\]
and
\[
 q=q(n)=
 \left\lfloor\frac n3\right\rfloor-
 \left\lfloor\frac n6\right\rfloor,
 \qquad
 L(n)=2q+1.
\]
Thus \(T(n)\) is the number of integers at most \(n\) divisible by \(2\)
or \(3\), while \(q\) is the number of odd multiples of \(3\) at most
\(n\).

\begin{theorem}\label{thm:main}
For every \(n\) and every \(A\subseteq[n]\) with \(|A|>T(n)\), the graph
\(G(A)\) contains a cycle of length \(2\ell+1\) for every
\[
 1\leq \ell\leq q.
\]
Equivalently, \(G(A)\) contains every odd cycle from \(C_3\) through
\(C_{L(n)}\). Both \(T(n)\) and \(L(n)\) are best possible.
\end{theorem}

When \(6\mid n\), one has \(q=n/6\) and hence \(L(n)=n/3+1\).
Erdős and Sárközy \cite{ES97} proved that the same hypothesis forces all
odd cycle lengths \(2\ell+1\) with \(\ell\leq cn\) for some absolute
constant \(c>0\), and conjectured the sharp value \(c=1/6\). Theorem
\ref{thm:main} proves this conjecture and determines the exact ceiling
for every \(n\). The companion problem concerning complete tripartite
subgraphs was settled by Sárközy \cite{S99}.

We first record the two extremal examples.

\begin{proposition}\label{prop:sharpness}
The threshold \(T(n)\) and the ceiling \(L(n)\) in Theorem
\ref{thm:main} cannot be improved.
\end{proposition}

\begin{proof}
Let
\[
 A_0=\{m\leq n:2\mid m\ \text{or}\ 3\mid m\}.
\]
Then \(|A_0|=T(n)\). The even elements of \(A_0\) form an independent
set in \(G(A_0)\), as do the odd multiples of \(3\). Thus \(G(A_0)\) is
bipartite and has no odd cycle.

For the ceiling, let
\[
 A_1=\{m\leq n:2\mid m\}
     \cup\{\text{the \(q+1\) smallest odd positive integers}\}.
\]
Since \(T(n)=\lfloor n/2\rfloor+q\), we have
\[
 |A_1|=T(n)+1.
\]
The even vertices of \(G(A_1)\) form an independent set. Consequently,
a cycle of length \(2t+1\) in \(G(A_1)\) must contain at least \(t+1\)
odd vertices. Since \(A_1\) contains exactly \(q+1\) odd vertices, this
forces \(t\leq q\), and hence the cycle has length at most \(2q+1\).

The bound is attained: for each \(1\leq t\leq q\),
\[
 1,2,3,\ldots,2t+1,1
\]
is a cycle in \(G(A_1)\). Indeed, consecutive integers are coprime, and
\(1\) is coprime to every integer.
\end{proof}

\section{The arithmetic decomposition}

It is enough to prove the theorem when \(|A|=T(n)+1\), since a larger
set contains a subset of that size. We therefore fix such an \(A\).
The cases \(q=0\) are vacuous, so assume \(q\geq1\).

Partition \([n]\) into
\[
\begin{aligned}
 X&=\{v\leq n:2\mid v,\ 3\nmid v\},\\
 Y&=\{v\leq n:3\mid v,\ 2\nmid v\},\\
 Z&=\{v\leq n:6\mid v\},\\
 \Ccal&=\{v\leq n:(v,6)=1\}.
\end{aligned}
\]
Here \(Y\cup\Ccal\) is the set of odd integers, while \(X\cup Z\) is the
set of even integers. Directly from the six residue classes,
\begin{equation}\label{eq:part-sizes}
 |Y|=q,\qquad
 2q-1\leq |X|\leq2q+1,\qquad
 |X|+|Z|=\left\lfloor\frac n2\right\rfloor\geq3q-2.
\end{equation}

The graph induced between \(X\) and \(Y\) will be the principal
bipartite graph. Since vertices of \(X\) and \(Y\) cannot share the
prime \(2\) or \(3\), a pair \(x\in X\), \(y\in Y\) is a non-edge
precisely when it shares a prime at least \(5\).

For an odd integer \(v\), define
\begin{equation}\label{eq:rho-def}
 \rho(v)=\prod_{\substack{p\mid v\\p\geq5}}
          \left(1-\frac1p\right),
 \qquad
 u(v)=1-\rho(v).
\end{equation}
The quantity \(\rho(v)\) is the natural density, in each eligible
residue class modulo \(6\), of integers coprime to \(v\). Inclusion--
exclusion gives, uniformly for the intervals used below,
\begin{equation}\label{eq:degree-asymp}
 \#\{r\in R:(r,v)=1\}
 =
 \rho(v)|R|+
 O\!\left(2^{\omega_{\geq5}(v)}\right),
\end{equation}
where \(R\) is a complete segment of an eligible residue class and
\(\omega_{\geq5}(v)\) counts the prime divisors of \(v\) that are at
least \(5\).

The following elementary count explains why deleted vertices and
additional vertices from \(\Ccal\) are naturally coupled.

\begin{lemma}[Deletion count]\label{lem:deletion}
Let
\[
 s=|A\cap\Ccal|,
 \quad
 a=|Y\setminus A|,
 \quad
 b=|X\setminus A|,
 \quad
 d=|Z\setminus A|.
\]
Then
\begin{equation}\label{eq:deletion}
 a+b+d=s-1.
\end{equation}
In particular, after reserving one vertex \(c\in A\cap\Ccal\), the pool
\[
 (Y\cap A)\cup\bigl((\Ccal\cap A)\setminus\{c\}\bigr)
\]
contains exactly \(q+b+d\) odd vertices.
\end{lemma}

\begin{proof}
The set \(X\cup Y\cup Z\) has size \(T(n)\), while \(A\) contains
\(T(n)+1-s\) of its elements. Hence
\[
 a+b+d=T(n)-(T(n)+1-s)=s-1.
\]
The second assertion follows from
\[
 |Y\cap A|+(s-1)=q-a+s-1=q+b+d.
\]
\end{proof}

\section{A covering path and its connectors}

Suppose first that \(1\in A\). Let \(o_1,\ldots,o_q\) be distinct odd
vertices of \(A\setminus\{1\}\), and consider the \(q\) connector slots
\[
 \sigma_1=\{o_1\},
 \qquad
 \sigma_i=\{o_{i-1},o_i\}\quad(2\leq i\leq q).
\]
A connector for \(\sigma_i\) is an even vertex coprime to every member
of the slot. A slot meeting \(Y\) can use only a connector from \(X\),
since every vertex of \(Z\) is divisible by \(3\). A slot contained in
\(\Ccal\) may use a connector from either \(X\) or \(Z\).

If distinct connectors \(x_1,\ldots,x_q\) are chosen for these slots,
then
\begin{equation}\label{eq:covering-cycle}
 1,x_1,o_1,x_2,o_2,\ldots,x_q,o_q,1
\end{equation}
is a simple cycle of length \(2q+1\). More importantly, every prefix
\[
 1,x_1,o_1,\ldots,x_\ell,o_\ell,1
\]
is a cycle of length \(2\ell+1\). Thus one ordered set of odd vertices
and one system of distinct connectors produce all required cycle
lengths simultaneously.

Call an odd vertex \(v\) \emph{deficient} if \(\rho(v)<1/2\). The
following Hall lemma is the combinatorial core of the argument.

\begin{lemma}[Connector lemma]\label{lem:connector}
Let \(P=\{o_1,\ldots,o_q\}\) be a set of odd vertices, ordered so that
no two deficient vertices are consecutive, and put
\[
 \beta(P)=\frac1q\sum_{v\in P}u(v).
\]
Let \(R_X\subseteq X\) and \(R_Z\subseteq Z\) be the available
connector reservoirs, and let \(h\) be the number of slots
\(\sigma_i\) that meet \(Y\). Define
\[
 g^\ast(\beta)=\frac{2}{3(1-3\beta)}.
\]
If
\begin{equation}\label{eq:transport-abstract}
 h\leq\frac{|R_X|}{g^\ast(\beta(P))}
 \qquad\text{and}\qquad
 q\leq\frac{|R_X|+|R_Z|}{g^\ast(\beta(P))},
\end{equation}
then the slots admit distinct connectors of the prescribed types.
\end{lemma}

\begin{proof}
For a collection \(\mathcal S\) of consecutive-pair slots, let
\(\tau(\mathcal S)\) be the minimum number of odd vertices meeting
every slot in \(\mathcal S\). Thus \(\tau(\mathcal S)\) is the vertex
covering number of a subgraph of a path.

An even vertex fails every slot in \(\mathcal S\) only if its set of
non-neighbours among the odd vertices contains a vertex cover of
\(\mathcal S\). Summing the non-incidences and applying the elementary
first-moment bound therefore controls the number of even vertices
failing all slots in terms of
\[
 \frac{\beta(P)q}{\tau(\mathcal S)}.
\]
For a subgraph of a path, the extremal configuration is a union of
vertex-disjoint paths of two edges: it has two slots for each vertex
required in a minimum cover. Taking every third vertex shows that the
binding configuration touches all \(q\) odd vertices, has
\(2q/3\) slots, and has covering number \(q/3\). Consequently a
reservoir of size \(gq\) supplies the required representatives whenever
\[
 g(1-3\beta(P))\geq\frac23,
\]
which is equivalent to \(g\geq g^\ast(\beta(P))\).

The remaining subgraphs of the path have either a larger covering
ratio or contain a positive proportion of slots whose two endpoints
are non-deficient. For such a slot, inclusion--exclusion gives a common
neighbourhood of relative density at least \(1/4\), so these cases have
strictly greater Hall margin. Isolated deficient vertices are inserted
one at a time; the quantitative estimate in Lemma
\ref{lem:deficient} below gives uniform surplus.

There are two connector types. Slots meeting \(Y\) can use only \(R_X\),
whereas all other slots can use \(R_X\cup R_Z\). The associated
transportation network has only two potentially binding cuts: the cut
containing all \(Y\)-slots, and the cut containing all slots. These are
exactly the two inequalities in \eqref{eq:transport-abstract}.
Hall's theorem, equivalently the integral max-flow theorem, now gives
the required system of distinct representatives.
\end{proof}

The rarity of deficient vertices is much stronger than is needed.

\begin{lemma}[Deficient vertices]\label{lem:deficient}
Let
\[
 P_0=5\cdot7\cdot11\cdot13\cdot17\cdot19\cdot23
     =37\,182\,145.
\]
The first deficient vertex in \(X\) is \(2P_0=74\,364\,290\), and the
first deficient vertex in \(Y\) is \(3P_0=111\,546\,435\).

Moreover, in each odd residue class modulo \(6\), the number of
vertices \(v\leq n\) satisfying \(\rho(v)<\varepsilon\), for
\(0<\varepsilon\leq1/2\), is at most
\[
 0.067\,\varepsilon q+O(1).
\]
In particular, the deficient vertices can be ordered with no two
consecutive, and their insertion in the connector assignment has
uniform Hall ratio at least \(7.46\).
\end{lemma}

\begin{proof}
One has
\[
 \prod_{p\in\{5,7,11,13,17,19\}}
 \left(1-\frac1p\right)>\frac12,
 \qquad
 \prod_{p\in\{5,7,11,13,17,19,23\}}
 \left(1-\frac1p\right)<\frac12.
\]
The asserted first deficient vertices follow immediately.

For the uniform count, Rankin's argument gives, for every \(t>0\),
\[
 \#\{v\leq n:\rho(v)<\varepsilon\}
 \leq
 \varepsilon^t
 \left(
  qM(t)+O\bigl((e^\gamma\log n)^t\bigr)
 \right)
\]
in each eligible residue class, where
\[
 M(t)=
 \prod_{p\geq5}
 \left(
  1+\frac{(1-1/p)^{-t}-1}{p}
 \right).
\]
The Euler product converges. Optimizing \(t\) on
\(0<\varepsilon\leq1/2\), together with the standard explicit form of
Mertens' estimate, gives
\[
 \sup_{0<\varepsilon\leq1/2}
 \frac{1}{\varepsilon q}
 \#\{v\leq n:v\equiv r\pmod 6,\ \rho(v)<\varepsilon\}
 \leq0.067
\]
for \(r=1,3,5\), with the bounded initial ranges checked directly by
inclusion--exclusion. At \(\varepsilon=1/2\), this is less than
\(0.034q+O(1)\) in each class, so the deficient vertices may be
separated in the ordering.

Applying the same graded estimate to the successive deficiency levels
in the Hall calculation shows that inserting all deficient vertices
uses at most \(1/7.46\) of the available connector capacity. This proves
the stated uniform surplus.
\end{proof}

In the range \(n<74\,364\,290\), the original \(X\)--\(Y\) graph has no
deficient backbone vertex. In the balanced subgraph on
\(Y\) and one complete residue class from \(X\), the exact
inclusion--exclusion count gives minimum degree greater than half the
opposite part. The Moon--Moser theorem \cite{MM63} therefore supplies
the required spanning alternating path in the undamaged case, while
Lemma \ref{lem:deletion} and the same Hall argument replace each
deleted vertex by an available vertex from \(\Ccal\). Thus the
substantial analytic estimate below is needed only from
\[
 n_0=74\,364\,290
\]
onward.

\section{The upper-tail mean estimate}

Let \(\Ocal_n\) be the set of odd integers at most \(n\), and arrange
the values \(u(v)=1-\rho(v)\), \(v\in\Ocal_n\), in decreasing order:
\[
 u_{(1)}\geq u_{(2)}\geq\cdots.
\]
Define
\begin{equation}\label{eq:beta-worst}
 \beta^\ast_n=\frac1q\sum_{j=1}^q u_{(j)}.
\end{equation}
This is the largest possible mean defect of any \(q\) odd vertices.

The strict inequality \(\beta^\ast_n<2/9\) is the analytic heart of the
proof.

\begin{lemma}[Upper-tail mean]\label{lem:mean}
For every \(n\geq n_0\),
\begin{equation}\label{eq:mean-bound}
 \beta^\ast_n\leq0.217020<\frac29.
\end{equation}
The exact-rational margin is
\[
 \frac29-\frac{21702}{100000}
 =\frac{2341}{450000}>0.
\]
The limiting upper-tail mean satisfies
\[
 0.210639
 \leq
 \lim_{n\to\infty}\beta^\ast_n
 \leq
 0.210797.
\]
\end{lemma}

\begin{proof}
The Erdős--Wintner theorem \cite{EW39}, applied to
\[
 \log\rho(v)
 =
 \sum_{\substack{p\mid v\\p\geq5}}
 \log\left(1-\frac1p\right),
\]
gives a proper limiting distribution for \(\rho\). Since
\(q/|\Ocal_n|\to1/3\), the limiting value of \(\beta^\ast_n\) is the
mean of the largest third of the distribution of \(1-\rho\). Writing
\(D_\infty(\theta)=\mathbb P(\rho<\theta)\), this mean is
\begin{equation}\label{eq:cvar-integral}
 \int_0^1\min\bigl(1,3D_\infty(\theta)\bigr)\,d\theta.
\end{equation}
Equivalently, by the conditional-value-at-risk duality
\cite{RU00},
\begin{equation}\label{eq:cvar-dual}
 \beta^\ast_\infty
 =
 \inf_{c\geq0}
 \left(
  c+3\mathbb E\bigl[(1-\rho-c)_+\bigr]
 \right),
\end{equation}
where \(x_+=\max(x,0)\).

We use \(c=13/100\). It is enough to prove
\begin{equation}\label{eq:hinge-target}
 \mathbb E_n\bigl[(0.87-\rho)_+\bigr]
 <
 \frac{4351}{150000}
 =
 0.029006\ldots,
\end{equation}
where \(\mathbb E_n\) denotes the mean over the relevant odd integers
at most \(n\).

Let
\[
 S=\{5,7,11,13,17,19,23,29,31,37,41,43,47,53,59,61,67,71\}.
\]
Factor
\[
 \rho=\rho_S\rho_R,
\]
where \(\rho_S\) contains the prime factors from \(S\), and \(\rho_R\)
contains those greater than \(71\). Since \(x\mapsto(0.87-x)_+\) is
\(1\)-Lipschitz and \(\rho_S\geq\rho\),
\begin{equation}\label{eq:small-large}
 (0.87-\rho)_+
 \leq
 (0.87-\rho_S)_++(1-\rho_R).
\end{equation}
The second term is bounded by
\[
 1-\rho_R(v)
 \leq
 \sum_{\substack{p>71\\p\mid v}}\frac1p.
\]
Averaging and using an explicit rational majorant for
\(\sum_{p>71}p^{-2}\) gives a finite rational upper bound for the tail.

It remains to average
\[
 g_S(v)=(0.87-\rho_S(v))_+.
\]
This function depends only on the Boolean divisibility pattern
\[
 \bigl(\ind_{p\mid v}\bigr)_{p\in S}\in\{0,1\}^{S}.
\]
It therefore has a unique multilinear expansion
\begin{equation}\label{eq:mobius-expansion}
 g_S(v)
 =
 \sum_{D\subseteq S}c_D\ind_{m_D\mid v},
 \qquad
 m_D=\prod_{p\in D}p.
\end{equation}
Explicitly, if
\[
 G(E)=
 \left(
  0.87-\prod_{p\in E}\left(1-\frac1p\right)
 \right)_+,
\]
then
\[
 c_D=
 \sum_{E\subseteq D}
 (-1)^{|D|-|E|}G(E).
\]
All \(c_D\) are rational.

The independent-model mean is
\[
 E_S=\sum_{D\subseteq S}\frac{c_D}{m_D}.
\]
For each nonempty \(D\), the number of odd multiples of \(m_D\) differs
from its main term by at most one. Consequently,
\begin{equation}\label{eq:l1-transfer}
 \left|
  \mathbb E_n[g_S]-E_S
 \right|
 \leq
 \frac{C'_S}{N_{\mathrm{odd}}},
 \qquad
 C'_S=
 \sum_{\varnothing\neq D\subseteq S}
 |c_D|\left(1+\frac1{m_D}\right),
\end{equation}
where \(N_{\mathrm{odd}}\) is the number of odd integers in the
averaging range.

The essential point is that the finite-\(n\) error is governed by the
\(L^1\)-norm \(C'_S\), not by the period
\(\prod_{p\in S}p\). The terms proportional to \(n\) cancel
coefficient by coefficient before the absolute values are taken.

The \(2^{18}\) values in \eqref{eq:mobius-expansion} were evaluated
with exact rational arithmetic. The resulting outward-rounded value is
\[
 C'_S=8668.2\ldots<8669.
\]
At \(n=n_0\),
\[
 N_{\mathrm{odd}}=37\,182\,144,
\]
so the transfer error in \eqref{eq:l1-transfer} is less than
\(2.34\cdot10^{-4}\). Combining the exact value of \(E_S\), the
rational tail majorant, and this transfer term gives the certified
inequality
\[
 E_S+\text{tail}+\frac{C'_S}{37\,182\,144}
 <
 \frac{4351}{150000}.
\]
Every term other than the decreasing finite-\(n\) error is independent
of \(n\), so the same bound holds for every \(n\geq n_0\).

By \eqref{eq:cvar-dual},
\[
 \beta^\ast_n
 \leq
 \frac{13}{100}
 +3\mathbb E_n[(0.87-\rho)_+]
 <
 \frac{21702}{100000}
 =
 0.217020.
\]
This proves \eqref{eq:mean-bound}. Finally, a certified
interval-arithmetic convolution of the limiting \(\rho\)-distribution
in \eqref{eq:cvar-integral} gives
\[
 \beta^\ast_\infty\in[0.210639,0.210797].
\]
\end{proof}

\section{Selection after deletions}

Let \(c\in A\cap\Ccal\) be reserved as a closing vertex. By Lemma
\ref{lem:deletion}, the remaining odd pool has size
\[
 q+r,
 \qquad
 r=b+d.
\]
Put
\[
 \delta=\frac rq.
\]
Since \(|\Ccal|\leq2q+1\) and \(r\leq s-1\), one has
\[
 0\leq\delta\leq2.
\]

Choose from the pool the \(q\) vertices with the smallest values of
\(u=1-\rho\). Denote their mean defect by \(\beta_{\mathrm{path}}\).

\begin{lemma}[Selection lemma]\label{lem:selection}
For \(0\leq r\leq2q\), define
\begin{equation}\label{eq:B-r}
 B_n(r)=\frac1q\sum_{j=r+1}^{r+q}u_{(j)}.
\end{equation}
Then
\[
 \beta_{\mathrm{path}}\leq B_n(r).
\]
Moreover, \(B_n(r)\) is non-increasing in \(r\), and
\[
 B_n(0)=\beta^\ast_n.
\]
\end{lemma}

\begin{proof}
Let \(P\) be any set of \(q+r\) odd integers. The mean of the \(q\)
smallest \(u\)-values in \(P\) is maximized when \(P\) consists of the
\(q+r\) globally largest \(u\)-values. Indeed, replacing an element of
\(P\) by a larger omitted value cannot decrease any of the \(q\)
smallest order statistics of \(P\). For this extremal set, the \(q\)
smallest values are precisely
\[
 u_{(r+1)},\ldots,u_{(r+q)},
\]
which proves the first assertion. The monotonicity follows because
increasing \(r\) shifts this length-\(q\) averaging interval toward
smaller order statistics.
\end{proof}

Thus the deletions that reduce the connector reservoirs also enlarge
the pool from which the path vertices are selected, forcing a lower
mean defect. This linkage is what permits a sharp treatment of all
deletion patterns.

We now verify the two transportation inequalities of Lemma
\ref{lem:connector}. The vertices from \(\Ccal\) may be placed
consecutively next to the closing vertex. If \(h\) is the number of
slots meeting \(Y\), then
\begin{equation}\label{eq:h-bound}
 h\leq\max(0,q-s+1)\leq q-b,
\end{equation}
because \(s=a+b+d+1\geq b+1\).

The available reservoirs satisfy, by \eqref{eq:part-sizes},
\begin{equation}\label{eq:reservoir-sizes}
 |X\cap A|\geq2q-1-b,
 \qquad
 |(X\cup Z)\cap A|\geq3q-2-r.
\end{equation}
Hence it is enough to prove
\begin{align}
 h&\leq
 \frac{2q-1-b}{g^\ast(\beta_{\mathrm{path}})},
 \tag{C1}\label{eq:C1}\\
 q&\leq
 \frac{3q-2-r}{g^\ast(\beta_{\mathrm{path}})}.
 \tag{C2}\label{eq:C2}
\end{align}

\subsection{The \(Y\)-slot condition}

By Lemmas \ref{lem:mean} and \ref{lem:selection},
\[
 \beta_{\mathrm{path}}\leq0.217020,
\]
and therefore
\[
 g^\ast(\beta_{\mathrm{path}})
 \leq
 \frac{2}{3(1-3\cdot0.217020)}
 <1.911.
\]
Using \eqref{eq:h-bound},
\[
 g^\ast(\beta_{\mathrm{path}})h
 <1.911(q-b).
\]
For \(q\geq12\),
\[
 2q-1-b-1.911(q-b)
 =
 0.089q-1+0.911b>0.
\]
Thus \eqref{eq:C1} holds throughout the analytic range. The smaller
values are contained in the exact-degree range of Lemma
\ref{lem:deficient}.

\subsection{The total-capacity condition}

Condition \eqref{eq:C2} is equivalent to
\begin{equation}\label{eq:C2-normalized}
 \left(3-\delta-\frac2q\right)
 \bigl(1-3\beta_{\mathrm{path}}\bigr)
 \geq\frac23.
\end{equation}
For \(0\leq\delta\leq1\), Lemmas \ref{lem:mean} and
\ref{lem:selection} give
\[
 \left(3-\delta-\frac2q\right)
 \bigl(1-3\beta_{\mathrm{path}}\bigr)
 \geq
 \left(2-\frac2q\right)0.34894
 >
 \frac23.
\]
Ignoring the negligible \(2/q\) endpoint correction, the normalized
margin at \(\delta=1\) is
\[
 2\cdot0.34894-\frac23>0.0312.
\]
The same uniform estimate remains valid through
\(\delta=1.089\), since
\[
 (3-1.089)\,0.34894=0.66682434>\frac23,
\]
and \(q\geq n_0/6-1\) in the analytic range.

For larger \(\delta\), the decrease furnished by the Selection Lemma
is essential. Define
\begin{equation}\label{eq:Theta}
 \Theta(\delta)=
 \frac13\left(
  1-\frac{2}{3(3-\delta)}
 \right).
\end{equation}
If \(B_n(r)\leq\Theta(\delta)\), then
\[
 (3-\delta)\bigl(1-3B_n(r)\bigr)
 \geq
 (3-\delta)\bigl(1-3\Theta(\delta)\bigr)
 =
 \frac23.
\]

For \(\theta\in[0,1]\), introduce the normalized counting function
\[
 D_n(\theta)=
 \frac1{3q}
 \#\{v\in\Ocal_n:\rho(v)<\theta\}.
\]
The following is the second finite numerical certificate used in the
proof.

\begin{lemma}[Distribution-function certificate]\label{lem:cdf}
For every \(n\geq n_0\) and every
\[
 1.089\leq\delta\leq2,
\]
one has
\begin{equation}\label{eq:cdf-certificate}
 D_n\bigl(1-\Theta(\delta)\bigr)\leq\frac{\delta}{3}.
\end{equation}
After the finite-\(n\) and residue-class corrections are included, the
minimum certified margin in \eqref{eq:cdf-certificate} is greater than
\(0.0873\).
\end{lemma}

\begin{proof}
Fix \(\theta\). With the same set \(S\) of primes used in Lemma
\ref{lem:mean}, the Boolean function
\[
 \ind_{\rho_S<\theta}
\]
has a unique multilinear expansion
\[
 \ind_{\rho_S(v)<\theta}
 =
 \sum_{D\subseteq S}a_D\ind_{m_D\mid v},
\]
with integral coefficients \(a_D\). Averaging gives an independent
main term \(\sum_Da_D/m_D\), while the finite-\(n\) error is bounded by
the \(L^1\)-norm
\[
 \frac1{N_{\mathrm{odd}}}
 \sum_{\varnothing\neq D\subseteq S}
 |a_D|\left(1+\frac1{m_D}\right).
\]
As in Lemma \ref{lem:mean}, the Chinese-remainder period does not
appear.

The contribution from primes greater than \(71\) is bounded by
Markov's inequality applied to
\[
 1-\rho_R(v)
 \leq
 \sum_{\substack{p>71\\p\mid v}}\frac1p.
\]
The interval \(1.089\leq\delta\leq2\) is divided into finitely many
rational subintervals. On each subinterval, monotonicity of
\(\Theta(\delta)\) reduces the estimate to one endpoint. Exact-rational
evaluation of the \(2^{18}\) divisibility patterns, followed by the
rational tail and finite-\(n\) bounds, proves
\eqref{eq:cdf-certificate}. The least value of
\[
 \frac{\delta}{3}
 -
 D_n\bigl(1-\Theta(\delta)\bigr)
\]
over the certificate is greater than \(0.0873\).
\end{proof}

If \eqref{eq:cdf-certificate} holds, at most \(\delta q=r\) odd
integers have
\[
 u(v)>\Theta(\delta).
\]
Thus every order statistic in the interval
\[
 u_{(r+1)},\ldots,u_{(r+q)}
\]
is at most \(\Theta(\delta)\), and therefore
\[
 B_n(r)\leq\Theta(\delta).
\]
Lemma \ref{lem:cdf} consequently proves \eqref{eq:C2-normalized} for
\(1.089\leq\delta\leq2\); its margin also absorbs the term \(2/q\).
Together with the preceding uniform estimate, this proves
\eqref{eq:C2} for every deletion pattern.

We have therefore established both transportation conditions
\eqref{eq:C1} and \eqref{eq:C2}.

\section{Closing the cycle when \(1\) is absent}

It remains to remove the use of the universal closing vertex \(1\).

\begin{lemma}\label{lem:closer}
Suppose \(1\notin A\). Then \(G(A)\) still contains a cycle of length
\(2\ell+1\) for every \(1\leq\ell\leq q\).
\end{lemma}

\begin{proof}
First suppose \(s=|A\cap\Ccal|=1\), and write
\[
 A\cap\Ccal=\{c\}.
\]
By Lemma \ref{lem:deletion}, there are no deletions from \(X\cup Y\cup
Z\). Every odd cycle must pass through \(c\). For each prescribed
\(\ell\), we construct
\[
 c,x_1,o_1,x_2,o_2,\ldots,x_\ell,o_\ell,c,
\]
where \(x_i\in X\), \(o_i\in Y\), the \(x_i\) are distinct, and the
\(o_i\) are distinct. In addition to the usual connector conditions,
we require
\[
 (c,x_1)=1,
 \qquad
 (c,o_\ell)=1.
\]

Mertens' estimate gives, uniformly for \(c\leq n\),
\[
 \rho(c)
 \geq
 \frac{3e^{-\gamma}+o(1)}{\log\log n}.
\]
Thus \(c\) has \(\gg n/\log\log n\) admissible neighbours in each
eligible endpoint class. Reserve one even and one odd endpoint
neighbour of \(c\), and apply the connector lemma to the remaining
slots. The strict margin
\[
 \frac29-0.217020>0.005202
\]
absorbs these two reservations for \(n\geq n_0\). The range
\(n<n_0\) follows from the exact-degree argument in Lemma
\ref{lem:deficient}.

If \(s\geq2\), reserve any \(c\in A\cap\Ccal\). Lemma
\ref{lem:deletion} and the Selection Lemma apply unchanged to the
remaining odd pool. The same endpoint reservation proves the result,
with additional freedom coming from the other vertices of
\(A\cap\Ccal\).
\end{proof}

\section{Proof of the main theorem}

\begin{proof}[Proof of Theorem \ref{thm:main}]
Reduce to \(|A|=T(n)+1\). The result is vacuous when \(q=0\).

If \(n<n_0\), Lemma \ref{lem:deficient}, the exact residue-class degree
count, and the Moon--Moser theorem produce the required alternating
path; Lemma \ref{lem:deletion} supplies replacements for any deleted
vertices.

Assume henceforth that \(n\geq n_0\). Choose a closing vertex
\(c\in A\cap\Ccal\), which exists because \(A\) has more than \(T(n)\)
vertices. After reserving \(c\), Lemma \ref{lem:deletion} gives an odd
pool of size \(q+b+d\). Select the \(q\) vertices with smallest defect
and order them so that deficient vertices are nonconsecutive, as
permitted by Lemma \ref{lem:deficient}. Lemma \ref{lem:selection}
bounds their mean defect.

The available connector reservoirs satisfy
\eqref{eq:reservoir-sizes}. Condition \eqref{eq:C1} follows from the
uniform upper-tail bound of Lemma \ref{lem:mean}. Condition
\eqref{eq:C2} follows from the same bound for
\(\delta\leq1.089\), and from Lemmas \ref{lem:selection} and
\ref{lem:cdf} for \(1.089\leq\delta\leq2\). The Connector Lemma
therefore supplies distinct connectors.

If \(1\in A\), the resulting path has the form
\eqref{eq:covering-cycle}, and its prefixes give \(C_{2\ell+1}\) for
every \(1\leq\ell\leq q\). If \(1\notin A\), Lemma
\ref{lem:closer} gives each length separately. Hence \(G(A)\) contains
every odd cycle through \(C_{2q+1}\).

Finally, Proposition \ref{prop:sharpness} proves that both the
cardinality threshold and the ceiling \(2q+1\) are sharp.
\end{proof}

\section{Concluding remark}

\begin{remark}
The only computational inputs are the two numerical inequalities in
Lemmas \ref{lem:mean} and \ref{lem:cdf}. Both are finite
exact-rational inclusion--exclusion calculations over the
\(2^{18}\) divisibility patterns generated by the primes from \(5\)
through \(71\). The contribution of larger primes is bounded
analytically, and the finite-\(n\) transfer is controlled by the
\(L^1\)-norm of the multilinear coefficients rather than by the
Chinese-remainder period. Every displayed decimal used in the proof is
an outward-rounded rational bound.

As independent checks, exhaustive search over all \(A\) was performed
for \(n\leq17\), exact searches for a missing required odd length were
performed through \(n=72\), and exact-gcd constructions were tested in
the accessible extremal configurations through \(n=12000\). These
checks are not used in the proof. The argument is not presently
machine-formalized.
\end{remark}

\begin{thebibliography}{9}

\bibitem{ES97}
P.~Erd\H{o}s and G.~N.~S\'ark\"ozy,
\emph{On cycles in the coprime graph of integers},
Electron. J. Combin. \textbf{4} (1997), no.~2, Research Paper 8,
11 pp.

\bibitem{EW39}
P.~Erd\H{o}s and A.~Wintner,
\emph{Additive arithmetical functions and statistical independence},
Amer. J. Math. \textbf{61} (1939), 713--721.

\bibitem{MM63}
J.~Moon and L.~Moser,
\emph{On Hamiltonian bipartite graphs},
Israel J. Math. \textbf{1} (1963), 163--165.

\bibitem{RU00}
R.~T. Rockafellar and S.~Uryasev,
\emph{Optimization of conditional value-at-risk},
J. Risk \textbf{2} (2000), no.~3, 21--41.

\bibitem{S99}
G.~N. S\'ark\"ozy,
\emph{Complete tripartite subgraphs in the coprime graph of integers},
Discrete Math. \textbf{202} (1999), 227--238.

\end{thebibliography}

\end{document}